% Sablon pentru realizarea lucrarii de licenta, conform cu recomandarile
% din ghidul de redactare:
% - https://fmi.unibuc.ro/finalizare-studii/

% Multumiri lui Gabriel Majeri, acest sablon a fost creat pe baza
% codului sursa a lucrarii sale de licenta. 
% Codul sursa: https://github.com/GabrielMajeri/bachelors-thesis
% Website: https://www.gabrielmajeri.ro/
%
% Acest sablon este licentiat sub Creative Commons Attribution 4.0 International License.
% modificat 8 noiembrie 2024
% !TEX program = xelatex
\documentclass[12pt, a4paper]{article}

% Suport pentru diacritice și alte simboluri
\usepackage{fontspec}

% Font Times New Roman
\setmainfont{Times New Roman} % This replaces \usepackage{times}

\usepackage{polyglossia}

% Suport pentru mai multe limbi

% Setează limba textului la română
\setdefaultlanguage{romanian}
% Am nevoie de engleză pentru rezumat
\setotherlanguages{english}

% Indentează și primul paragraf al fiecărei noi secțiuni
\SetLanguageKeys{romanian}{indentfirst=true}

% Suport pentru diferite stiluri de ghilimele
\usepackage{csquotes}

\DeclareQuoteStyle{romanian}
  {\quotedblbase}
  {\textquotedblright}
  {\guillemotleft}
  {\guillemotright}

% Utilizează biblatex pentru referințe bibliografice
\usepackage[
    style=ieee,
    backend=biber,
    maxbibnames=50,
    sorting=none
]{biblatex}

\addbibresource{bibliography.bib}

% Setează spațiere inter-linie (doublespacing din LaTeX este echivalentul setarii 1.5 in Microsoft Word)
\usepackage{setspace}
\doublespacing

% Modificarea geometriei paginii (cu margini de 2,5 cm) 
\usepackage[margin=2.5cm]{geometry}

% Include funcțiile de grafică
\usepackage{graphicx}
% Încarcă imaginile din directorul `images`
\graphicspath{{./images/}}

% Listări de cod
\usepackage{listings}

% Linkuri interactive în PDF
\usepackage[
    colorlinks,
    linkcolor={black},
    menucolor={black},
    citecolor={black},
    urlcolor={blue}
]{hyperref}

% Simboluri matematice codificate Unicode

% Comenzi matematice
\usepackage{amsmath}
\usepackage{mathtools}
\usepackage[warnings-off={mathtools-colon,mathtools-overbracket}]{unicode-math}

% Formule matematice
\newcommand{\bigO}[1]{\symcal{O}\left(#1\right)}
\DeclarePairedDelimiter\abs{\lvert}{\rvert}

% Suport pentru rezumat în două limbi
% Bazat pe https://tex.stackexchange.com/a/70818
\newenvironment{abstractpage}
  {\cleardoublepage\vspace*{\fill}\thispagestyle{empty}}
  {\vfill\cleardoublepage}
\renewenvironment{abstract}[1]
  {\bigskip\selectlanguage{#1}%
   \begin{center}\bfseries\abstractname\end{center}}
  {\par\bigskip}

% Suport pentru anexe
\usepackage{appendix}

% Stiluri diferite de headere și footere
\usepackage{fancyhdr}

% Metadate
\title{Recunoasterea emotiilor printr-o abordare multimodala}
\author{Ionascu Andrei}

% Generează variabilele cu @
\makeatletter

\begin{document}

% --- PAGINA DE TITLU ---
\begin{titlepage}
\makeatletter % Permite folosirea comenzilor cu @

\newgeometry{left=2cm,right=2cm,bottom=1cm,top=2cm}

\begin{figure}[!htb]
    \centering
    \begin{minipage}{0.2\textwidth}
        % Asigură-te că fișierul există în folderul images/
        \includegraphics[width=\linewidth]{images/logo-ub.png} 
    \end{minipage}
    \begin{minipage}{0.5\textwidth}
        \large
        \begin{center}
            \textbf{UNIVERSITATEA DIN BUCUREȘTI} \\
            \vspace{0.3cm}
            \textbf{FACULTATEA DE \\ MATEMATICĂ ȘI INFORMATICĂ}
        \end{center}
    \end{minipage}
    \begin{minipage}{0.2\textwidth}
        \includegraphics[width=\linewidth]{images/logo-fmi.png}
    \end{minipage}
\end{figure}

\vspace{0.5cm}
\begin{center}
\textbf{SPECIALIZAREA INFORMATICĂ}
\end{center}

\vspace{2cm}

\begin{center}
\Large \textbf{Lucrare de licență}
\end{center}

\vspace{1cm}

\begin{center}
\huge \textbf{\MakeUppercase{\@title}}
\end{center}

\vfill % Împinge conținutul în jos

\begin{flushleft}
\large \textbf{Coordonator științific:} \\
Conf. Univ. Dr. Alexe Dumitru-Bogdan
\end{flushleft}

\vspace{-3.1cm} % Ajustare poziție text dreapta

\begin{flushright}
\large \textbf{Absolvent:} \\
\@author
\end{flushright}

\vspace{2cm}

\begin{center}
\Large \textbf{București, iunie 2026}
\end{center}

\makeatother % Închide accesul la comenzile cu @
\end{titlepage}

\restoregeometry % Revine la marginile normale
\newpage
\pagenumbering{roman} % Numerotare cu cifre romane pentru început

\subsection*{Abstract}
Pe masura ce inteligenta artificiala devine o componenta din ce in ce mai indispensabila in aplicatiile pe care oamenii le folosesc in fiecare zi, nevoia de sisteme care nu sunt doar precise, dar si eficiente din punct de vedere energetic si transparente cand vine vorba de deciziile pe care le iau a crescut considerabil. Aceasta lucrare exporeaza problema identificarii emotiilor umane printr-o abordare multimodala, analizand replici text, prozodie verbala si trasaturi faciale din setul de date MELD (Multimodal EmotionLines Dataset). Cercetarea exploreaza performantele obtinute de solutiile care integreaza un singur format de intrare, comparand ulterior imbunatatirile obtinute prin diverse metode de fuziune ale diferitelor surse de informatii.
Spre deosebire de abordarile traditionale care prioritizeaza performanta, aceasta teza pune accentul pe compromisul dintre acuratete si eficienta energetica si de timp, in timp ce transparentizeaza modul in care diferitele modele iau decizii, prin metode de XAI (Explainable AI)

\newpage

% --- INTRODUCERE ---

\section{Introducere}

\subsection{Context si evolutie istorica}
Recunoasterea emotiilor a evoluat de-a lungul timpului, plecand de la fluxuri conventionale, bazate pe caracteristici extrase manual pentru fiecare modalitate in parte, catre arhitecturi de tip deep learning, adesea pre-antrenate, capabile sa modeleze contextul, structura temporala si interactiunile cross-modale, dupa cum detaliaza \cite{PORIA201798}. Pe partea de audio si vorbire, domeniul a migrat de la descriptori clasici si clasificatori simpli catre modelele CNN-LSTM aplicate semnalului brut (raw waveform) \cite{Trigeorgis2016adieu} si, ulterior, catre codificatori de vorbire self-supervised, precum wave2vec 2.0 \cite{baevski2020wav2vec}. In vederea artificiala, traiectoria porneste de la descriptori de textura, precum Local Binary Patterns\cite{shan2009facial}, si ajunge la retele convolutionale, retele bazate pe mecanisme de atentie\cite{lideng2020deep} si modele FER optimizate pentru dispozitive de tip embedded\cite{tan2019efficientnet}. Procesarea limbajului natural cu scopul recunoasterii emotiilor a evoluat de la sisteme bazate pe lexicon si seturi de date restranse\cite{mohammad2013crowdsourcing}, la preantrenarea slab supravegheata si fine-tuning a modelelor de tip Transformer \cite{devlin2018bert} pe corpuri masive de date, cum este GoEmotions\cite{demszky-etal-2020-goemotions}. Pe langa modalitatile clasice, studiile recente exploreaza ramura de informatica fiziologica, unde lucrarile fundamentale s-au concentrat pe semnale EEG (Electroencefalograma) si periferice, utilizand seturi de date de referinta precum DEAP\cite{Koelstra2012DEAP}, progresand recent catre modele grafice care analizeaza topologia semnalelor cerebrale\cite{Song2020EEGER}. 

\subsection{Motivatie}
Motivatia acestei cercetari vine din faptul ca, odata cu integrarea functionalitatilor bazate pe inteligenta artificiala in tot mai multe produse software folosite in mod uzual de o tot mai mare parte din oameni, apare si nevoia de sisteme nu doar precise, dar si eficiente si de incredere.

\subsection{Problema: Transparenta si eficienta}

\subsubsection{Problema "Cutiei negre"}
Vasta majoritate a publicatiilor stiintifice din domeniul recunoasterii emotiilor au ca singur scop obtinerea unei precizii cat mai mari. Totusi, complexitatea necesara maximizarii acuratetii duce adesea la modele de tip "black-box", al caror proces de gandire este prea complicat pentru a fi inteles chiar de catre oameni de stiinta, cu atat mai putin de catre utilizatorul obisnuit. Acest lucru duce la o lipsa de incredere in progresul tehnologic si la o reticenta in folosirea noilor produse.

\subsubsection{Costul Energetic}
Complexitatea modelelor State-Of-The-Art presupune, in majoritatea cazurilor, si un cost energetic proportional. Dincolo de amprenta ecologica, nevoia masiva de resurse computationale face imposibila executia acestor modele inafara unor centre de date specializate, fapt ce limiteaza drastic portabilitatea si ridica provocari privind confidentialitatea datelor. In plus, dezvoltarea acestor solutii implica bariere financiare semnificative, cauzate nu doar de infrastructura necesara procesarii, dar si de costurile logistice aferente stocarii si gestionarii unor seturi de date voluminoase.

\subsection{Contributii si obiective}

\subsubsection{Integrarea multimodalitatii}
TO BE CONTINUED ONCE FULLY IMPLEMENTED

\subsubsection{Explicabilitatea rezultatelor}
TO BE CONTINUED ONCE FULLY IMPLEMENTED

\subsubsection{Compromisul dintre acuratete si eficienta energetica}
TO BE CONTINUED ONCE FULLY IMPLEMENTED
\newpage

\section{Fundamentare teoretica si stadiul actual al cercetarii}

\subsection{Recunoasterea automata a emotiilor (FER/SER/MER)}
Recunoașterea automată a emoțiilor reprezintă un subdomeniu al învățării automate afective care urmărește clasificarea sau estimarea stărilor emoționale pe baza unor semnale observabile. Scopul principal al acestui efort este de a crea sisteme si aplicatii care nu doar rezolva o problema, dar o fac interactionand cu utilizatorii tinand cont de anumite indicii nonverbale si paraverbale care in interactiunile pur umane sunt interpretate implicit. Potrivit lui Rosalind Picard, daca dorim calculatoare cu adevarat inteligente, capabile de interactiuni naturale, trebuie sa le oferim abilitatea de a recunoaste, intelege, poate chiar de a exprima emotii. \cite{picard1997affective}

În literatura de specialitate, emoțiile pot fi reprezentate fie prin categorii discrete, precum bucurie, tristețe, furie, frică, dezgust, surpriză sau neutralitate\cite{ekman1992argument}, fie prin modele dimensionale, în care starea afectivă este descrisă prin axe continue, precum valența și nivelul de activare\cite{russell1980circumplex}. În contextul acestei lucrări, problema este tratată ca o sarcină de clasificare multi-clasă, în care fiecare eșantion este asociat cu o etichetă emoțională discretă.

Recunoașterea emoțiilor poate fi formulată ca o problemă de clasificare supravegheată sau învățare de reprezentări, abordată fie unimodal, fie multimodal. În setările unimodale, informația emoțională este extrasă dintr-o singură sursă: text (analiză semantică și contextuală cu modele de tip Transformer), imagini faciale (FER), unde se utilizează rețele convoluționale sau arhitecturi bazate pe atenție pentru a învăța reprezentări spațiale robuste la variații de iluminare, poziție și ocluzii, sau semnal audio (SER), unde caracteristicile acustice (prosodie, energie, coeficienți cepstrali, spectrograme) sunt modelate prin CNN-uri, RNN-uri sau encodere self-supervised. Abordările tradiționale bazate pe feature engineering au fost în mare parte înlocuite de modele end-to-end care optimizează direct reprezentările latente. În contrast, recunoașterea multimodală (MER) integrează aceste surse eterogene prin strategii de fuziune (early, late sau hybrid fusion), capturând corelațiile inter-modale și reducând ambiguitatea semantică \cite{baltrusaitis2018multimodal}, aspect critic în contexte conversaționale complexe precum MELD\cite{poria2019meld}, unde emoția depinde de dinamica sincronizată dintre text, semnal acustic și expresii faciale.

\subsection{Modalitati de intrare si procesarea semnalelor}
În recunoașterea automată a emoțiilor, alegerea modalităților de intrare (text, audio, vizual) determină tipul de informație disponibil: textul captează conținutul semantic, audio-ul reflectă prosodia și modul de rostire, iar imaginile faciale surprind expresivitatea. Aceste modalități au structuri diferite (discretă secvențială, continuă temporală, respectiv spațială/spațio-temporală), necesitând metode distincte de preprocesare și transformare în reprezentări numerice. În sistemele multimodale, este esențială alinierea între modalități la nivel de instanță (ex. aceeași replică), pentru a permite fuziunea corectă a informației și a evita introducerea de zgomot. În această lucrare, textul, audio-ul și datele faciale sunt procesate inițial separat, pentru a evalua contribuția fiecărei modalități înainte de integrarea multimodală.

\subsubsection{Procesarea limbajului natural}
Procesarea limbajului natural transformă replicile textuale în reprezentări numerice utilizabile de modele de învățare automată. În recunoașterea emoțiilor, textul oferă informație semantică relevantă, deoarece anumite expresii pot semnala direct stări afective. Totuși, interpretarea emoției nu este întotdeauna trivială, deoarece sensul depinde de context (ex. ironie sau ambiguitate), motiv pentru care metodele moderne modelează relațiile dintre tokeni, nu doar prezența unor cuvinte-cheie.

Preprocesarea textuală include pași standard precum curățarea, normalizarea și tokenizarea, urmați de transformarea în vectori numerici. Metodele clasice (bag-of-words, TF-IDF, Word2Vec, GloVe) sunt simple, dar limitate contextual. În schimb, modelele bazate pe Transformer generează embeddings contextuale, unde reprezentarea unui cuvânt depinde de întreaga secvență, permițând captarea unor nuanțe semantice mai fine.\cite{vaswani2017attention}

În cadrul sistemelor multimodale, textul este o modalitate eficientă și compactă, dar incompletă, deoarece nu surprinde indicii non-verbale precum tonul sau expresivitatea facială. Prin urmare, performanța sa trebuie evaluată în raport cu modalitățile audio și vizuale.

\subsubsection{Analiza prozodica}
Analiza prozodică extrage informații emoționale din proprietățile acustice ale vorbirii, spre deosebire de NLP care se bazează pe conținutul semantic. Caracteristici precum frecvența fundamentală, energia, intensitatea, ritmul sau intonația pot semnala stări afective chiar și în absența unor indicii explicite în text. Procesarea audio implică segmentarea semnalului, normalizarea și împărțirea în ferestre temporale, urmate de extragerea unor reprezentări precum MFCC sau spectrograme Mel. MFCC-urile oferă o descriere compactă a spectrului, în timp ce spectrogramele păstrează structura timp-frecvență, fiind potrivite pentru modele neuronale.\cite{schuller2018speech}

O limitare majoră este variabilitatea ridicată a semnalului (diferențe între vorbitori, zgomot, condiții de înregistrare), ceea ce face ca indicii acustici să fie mai puțin stabili. Prin urmare, informația prozodică este utilă, dar trebuie tratată ca o sursă complementară în recunoașterea emoțiilor.

\subsubsection{Analiza faciala}
Analiza facială utilizează informația vizuală din regiunea feței pentru identificarea expresiilor asociate stărilor emoționale, fiind o componentă importantă a recunoașterii emoțiilor non-verbale. Mișcările musculare din zona ochilor, sprâncenelor și gurii oferă indicii relevante despre starea afectivă, iar problema este tratată, în mod obișnuit, ca o sarcină de analiză de imagini sau secvențe video.

Pipeline-ul începe cu extragerea cadrelor din videoclip, urmată de eșantionare pentru a reduce costul computațional. Se aplică apoi detecția feței, decuparea și normalizarea imaginilor, astfel încât modelul să se concentreze pe regiunea relevantă. În metodele clasice, expresiile erau descrise prin caracteristici extrase manual (LBP, muchii, landmarks), însă acestea sunt limitate în captarea variațiilor complexe. Modelele moderne bazate pe CNN învață automat reprezentări ierarhice și sunt mai robuste la variații de iluminare, poziție sau ocluzii.

În cazul datelor video, o singură imagine nu este întotdeauna suficientă, deoarece expresiile sunt dinamice. Deși modelarea secvențelor complete este costisitoare, o soluție practică constă în selectarea mai multor cadre relevante și agregarea informației la nivel de caracteristici sau predicții.

Totuși, această modalitate rămâne sensibilă la factori precum calitatea imaginii, poziția capului sau expresiile controlate voluntar. Din acest motiv, analiza facială este, în general, utilizată împreună cu alte modalități (text și audio), contribuind cu indicii vizuale complementare în cadrul sistemelor multimodale.

\subsection{Arhitecturi de tip Deep Learning pentru date secventiale}
Având în vedere că toate cele trei modalități analizate conțin o componentă secvențială — textul ca secvență de tokeni, audio-ul ca semnal temporal, iar videoclipul ca succesiune de cadre — alegerea arhitecturii de modelare devine esențială. Următorul subcapitol prezintă principalele arhitecturi de tip deep learning utilizate pentru procesarea datelor secvențiale și relevanța acestora în contextul recunoașterii emoțiilor.

\subsection{Strategii de fuziune multimodala}
Datele utilizate în recunoașterea emoțiilor sunt în mod natural secvențiale, ceea ce impune utilizarea unor modele capabile să captureze atât caracteristici locale, cât și dependențe temporale. Textul depinde de ordinea tokenilor, audio-ul de variațiile în timp ale semnalului, iar expresiile faciale pot evolua între cadre video. Din acest motiv, arhitecturile de deep learning sunt alese în funcție de capacitatea lor de a modela aceste dependențe.

Rețelele neuronale recurente (RNN), inclusiv variantele LSTM și GRU, procesează secvențele pas cu pas și mențin o stare internă care permite captarea contextului. LSTM-urile și GRU-urile rezolvă parțial problema dependențelor pe termen lung prin mecanisme de tip „gating”, fiind utilizate frecvent pentru text, audio sau secvențe video. Totuși, ele sunt mai greu de paralelizat și au costuri de antrenare mai ridicate.

Rețelele convoluționale (CNN) pot fi aplicate și pe date secvențiale, extrăgând tipare locale (ex. n-grame în text sau structuri timp-frecvență în audio). Avantajul lor este eficiența computațională și paralelizarea, dar modelarea dependențelor pe termen lung este limitată fără mecanisme suplimentare.

Arhitecturile bazate pe atenție, în special Transformer-ele, permit modelarea directă a relațiilor globale prin mecanisme de self-attention. Acestea sunt dominante în NLP (ex. BERT) și din ce în ce mai utilizate în audio și vizual, inclusiv în combinație cu preantrenare self-supervised (ex. wav2vec 2.0). Ele oferă reprezentări contextuale puternice, dar vin cu costuri computaționale ridicate.

În sisteme multimodale, alegerea arhitecturii trebuie să echilibreze performanța, costul și interpretabilitatea. O abordare modulară, în care fiecare modalitate este procesată separat (cu modele adaptate tipului de date) înainte de fuziune, permite o analiză mai clară a contribuției fiecărei surse și o integrare mai controlată.

\subsection{Inteligenta artificiala explicabila (XAI)}

\subsection{Eficienta energetica si computationala in AI}

\subsection{Setul de date MELD si stadiul actual al domeniului}
\newpage
\section{Metodologie}

\subsection{Analiza setului de date MELD}

\subsection{Preprocesarea datelor}

\subsection{Arhitecturi si modele propuse}

\subsection{Integrarea componentelor de explicabilitate}

\subsection{Strategia de masurare a eficientei}
\newpage
\section{Implementarea sistemelor}

\subsection{Tehnologii utilizate}

\subsection{Detalii de implementare}

\subsection{Antrenarea modelelor}
\newpage
\section{Rezultate experimentale si evaluare}

\subsection{Performanta bruta}

\subsection{Analiza XAI}

\subsection{Analiza eficientei}

\subsection{Discutii}
\newpage
\section{Concluzii si directii viitoare}


% --- CUPRINS ---
\tableofcontents
\newpage
\pagenumbering{arabic} % Revenire la cifre arabe pentru capitole

% --- BIBLIOGRAFIE ---
\renewcommand{\bibname}{References}
\printbibliography[title={Bibliografie}]

\end{document}